\section{Introduction}
\label{sec:introduction}

Legged robots enable mobility in
environments that are inaccessible to wheeled and tracked systems.
Their ability to traverse uneven, unstructured, and unpredictable
terrains, such as disaster zones, staircases, dense
vegetation, and rocky landscapes, positions them as key
enablers for future exploration, inspection, and rescue missions.
Unlike wheeled rovers, which are fundamentally limited by the continuous
support surface required for traction, quadrupedal systems can effectively
decouple their path of motion from the continuous topology of the ground,
selecting only the most viable, discrete footholds across otherwise
impassable terrain.

Realizing this potential requires a robot to continuously and
intelligently decide where and when to place its feet, a
process known as \textit{contact planning}. Existing approaches to
this problem broadly fall into three categories: traditional
optimization control, end-to-end learned control, and hybrid
methods that combine the two. Traditional optimization offers strong
guarantees but scales poorly once discrete contact decisions enter
the problem. End-to-end learning solves the mixture of discrete and
continuous decisions implicitly, but sacrifices interpretability and
formal safety. Hybrid methods aim to keep a learned component narrow,
selecting footholds or gait parameters, while a controller based on a
model retains responsibility for stability and tracking. Even
within hybrid designs, most existing planners still commit to a
fixed gait schedule. Contact timing is generated by a periodic phase
schedule or a library of gait primitives, and the learned component
only chooses \textit{where} to place a foot, not whether the
schedule itself should be broken. This constrains a robot's ability
to exploit the full versatility legged locomotion affords on terrain
that a schedule cannot anticipate.

\subsection{Terminology}
\label{subsec:introduction-terminology}

Legged robotics uses several overlapping terms inconsistently across
the literature; we fix their meaning here and use them consistently
throughout. \textit{Footstep planning} refers to the process of
selecting when and where to step. A planner is \textit{greedy} if it
provides only instantaneous actions, with no explicit lookahead over
future footsteps. \textit{Dynamic} locomotion is the opposite of
quasi-static locomotion. Balance is maintained through active control
and dynamic interaction with the environment, rather than by keeping
the center of mass within the support polygon at all times. In this
paper, \textit{dynamic} never means fast or reactive; it refers
specifically to this balance strategy. A \textit{gaited} or
\textit{gait-based} contact sequence is one generated by a periodic
phase schedule or by transitioning between a fixed library of gait
primitives; we call the opposite an \textit{acyclic} gait, a footstep
plan that cannot be represented as such phase cycles or primitive
transitions. \textit{Non-gaited} and \textit{free-gait} are used
synonymously with acyclic elsewhere in the literature; we use
\textit{acyclic} throughout for consistency.

\subsection{Related Work}

\subsubsection{Traditional Control Approaches}
Traditional quadruped locomotion commonly relies on explicit
optimization of footstep positions, contact sequences, and timing via
Model Predictive Control (MPC) \cite{Kim2019highly-dynamic}, which
handles continuous quantities such as joint torques and ground
reaction forces well \cite{Wensing2022Nov} but becomes a
Mixed-Integer Quadratic Program once discrete contact states are
introduced \cite{Geisert2019contact-planning}. Deits and Tedrake
\cite{deits_footstep_2014} formulated footstep placement as an MIQP
over convex terrain regions but restricted motion to sequential
placements without optimizing contact timing; Winkler et al.
\cite{winkler_gait_2018} and Aceituno-Cabezas et al.
\cite{aceituno_simultaneous_2019} extended this to jointly optimize
contact scheduling and timing, at a computational cost that precludes
reactive control in real time. Alternatives based on search, such as
Monte Carlo Tree Search \cite{taouil_non-gaited_2024,
akizhanov_learning_2024}, generate genuinely acyclic sequences, but
high branching factors keep them too slow for correction at high
frequency.

\subsubsection{End-to-End Neural Networks}
End-to-end learning methods map sensor data directly to joint
targets, bypassing explicit dynamics modeling and inherently solving
the mixed-integer contact problem \cite{zhang_learning_2023,
zhang_novel_2024}. This flexibility comes at the cost of formal
stability guarantees, large training data requirements, and
unreliable generalization outside the training distribution.

\subsubsection{Hybrid Methods and Learned Gait Policies}
Hybrid architectures merge learning components with execution based
on a model to balance tractability and stability. Several methods
focus on gait modulation at a high level. Shi et al.
\cite{shi_terrain-aware_2023} modulate trajectory parameters for
efficiency, Xie et al. \cite{xie_glide_2023} train RL policies on
centroidal dynamics given fixed gait heuristics, and blind
proprioceptive approaches \cite{duan_sim--real_2022,
siekmann_blind_2021, lee_learning_2020} push robust traversal without
exteroception. Other systems explicitly select among gaits via Deep
Q-Networks \cite{da_learning_2020}, neural phase integrators
\cite{yang_fast_2021}, or offline gait libraries integrated into
MPC running at high frequency \cite{sun_online_2024}. In every case, the learned
component operates on top of a periodic gait representation and
cannot itself produce contact sequences outside it.

\subsubsection{Learning-Based Footstep Planners}
Closer to our setting, explicit footstep planners driven by data
couple decisions about footholds with an MPC solver. Early architectures such
as FootstepNet \cite{gaspard_footstepnet_2024}, DeepLoco
\cite{peng_deeploco_2017}, and DeepGait \cite{tsounis_deepgait_2020}
laid the groundwork but localized to bipeds or fixed cyclic stepping.
RLOC \cite{gangapurwala_rloc_2022}, Omar et al.
\cite{omar_fast_2022}, and Villarreal et al.
\cite{villarreal_fast_2019} leverage convolutional terrain inputs
combined with fixed or periodic timing. Learned steppability terrain
maps \cite{asselmeier_steppability-informed_2024,
omar_safesteps_2023} inform where to place a foot but not when,
leaving contact timing to an external schedule. Most directly
relevant, ContactNet \cite{bratta_contactnet_2024} uses a
convolutional evaluator to extract non-gaited footstep costs from a
discretized grid, but its selection procedure is restricted to
moving one leg at a time in strict sequence, which caps how far a
system built on it can depart from behavior resembling a gait even
when the underlying cost estimates are acyclic.

\subsection{Motivation and Contributions}
Despite this progress, a clear gap remains. Hybrid planners either
inherit a periodic gait schedule from the control layer below them,
or, like ContactNet, remove the schedule but still commit footsteps
one leg at a time by construction. This motivates our core question:
whether a planner that commits to at most one footstep at a time,
re-deciding faster than a swing completes, can produce dynamic
acyclic gaits that a scheduled or sequential planner cannot.

GaitNet answers this by never representing coordination between legs
explicitly. At each control tick it evaluates a small set of
footstep candidates filtered by terrain per leg, locally refines them
against its own learned value function, and selects at most one
action, or none, for the entire robot. Because a selected swing
outlasts many such ticks, the planner can commit a second leg while
the first is still airborne. Any coordination between legs that
results is an emergent consequence of re-deciding quickly, not
something the architecture was given to reason about, though the
policy may learn to anticipate it implicitly through reinforcement
learning.

Our primary contributions are the following.
\begin{enumerate}
  \item A greedy footstep planner that emits at most one action per
    control tick, from which swing across multiple legs emerges
    without any scheduled phase representation, together with a
    simplified action space that removes the need for the footstep
    evaluation network with two stages used in prior work.
  \item A single-network formulation in which a trained value
    function is reused at deployment time as a local refinement
    objective over a small set of candidates filtered by terrain.
  \item Validation on a Unitree Go1 quadruped against the Legged
    Perceptive Stack, a leading perceptive Nonlinear Model Predictive
    Control (NMPC) baseline, together with ablations of GaitNet
    itself with restricted concurrency that isolate the contribution
    of emergent swing across legs.
\end{enumerate}

This paper builds on our preliminary framework, GaitNet, introduced in
\cite{Sullivan2025}, which established the single-footstep,
single-tick action space and reported promising results in simulation
under a discrete, randomly sampled candidate distribution. Here we
extend that formulation with a continuous candidate search refined by
gradients in place of discrete resampling, and provide the first
hardware validation of the resulting system. Simulation results
attributed to the preliminary framework are reported for context in
\autoref{sec:results}. We have not rerun those experiments under the
system described here, and label such results accordingly throughout.
